\section{What We Also Tried}
\label{sec:also-tried}

The objective in Eq.~\ref{eq:objective} has three terms.
An earlier version of this method had seven, listed in
Table~\ref{tab:retired}.
This appendix records what those terms were, why each seemed justified, and
what happened when they were measured in a regime where they could have been
right.
We report them because the reduction was the most informative result of the
project, and because a reader deciding whether to build a similar objective
deserves the negative half of our evidence.

\begin{table}[t]
\caption{The earlier seven-term objective. ``Weight'' is the shipped value in
that configuration. Only the first two terms and the M7 function anchor remain.
Section references point to the measurement that retired each term.}
\label{tab:retired}
\centering
\small
\begin{tabular}{llll}
\toprule
Term & Weight & Status & See \\
\midrule
Soft cross-entropy & 1.0 & retained & Sec.~\ref{sec:crossres} \\
Soft Dice & 0.25 & retained, reweighted & Sec.~\ref{sec:crossres} \\
M7 function anchor & 0.5 & retained & Sec.~\ref{sec:crossres} \\
Medial crest recall & 1.0 & retired & \ref{sec:also-crest} \\
Separation shell & 2.0 & retired & \ref{sec:also-separation} \\
M7 preservation & 1.0 & retired & \ref{sec:also-preservation} \\
Dynamic connectivity & 0.03125 & retired & \ref{sec:also-connectivity} \\
\midrule
Trust region & --- & retired & App.~\ref{sec:trust} \\
\bottomrule
\end{tabular}
\end{table}

\subsection{Why the Stack Survived as Long as It Did}
\label{sec:also-regime}

The terms were not retained because they had been shown to work.
They were retained because, in the regime where they were developed, nothing
could be shown either way, and a term that shows no measurable effect is easy to
keep.
Three properties of that regime were jointly responsible.

First, the parameter-space trust region was saturated: the post-update
projection was active on essentially every optimizer step, so no term could
move the function far enough to demonstrate either benefit or harm.
Appendix~\ref{sec:trust} treats this in detail because it is the load-bearing
failure.

Second, the promotion gate rewarded the behaviour the terms produced.
Its dominant criterion required that at least 95\% of the M7 voxels already
judged teacher-correct survive within a one-voxel neighbourhood of the
student's prediction.
That is a retention test, and preservation and the anchor are retention terms.
A thinner, better student fails it.
Cross-referencing the artifacts of that era, no shipped term scored above the
no-term baseline on this gate; the only intervention that moved it was an
inference-time blend with M7, which was itself later withdrawn as confounded.
A companion anti-blob check compounded the problem by pinning its reference to
an older prediction grid with a foreground cap, which bounded any candidate at
roughly 39\% of the foreground of a previously accepted model and made growth
scoreable only as a violation.

Third, the training corpus had narrowed to 4,096 rows from a single scroll.
A shape prior is most attractive exactly when data is scarce, and least
testable for the same reason.

A fourth, procedural point belongs here.
The terms were added incrementally and evaluated as a stack against the previous
stack, so what was tuned was the interaction, not the terms.
Leave-one-out arms were specified for five of them and never trained.
Preservation is the clearest casualty of this ordering: it was always paired
with the separation shell that offsets its cost, so the question of what it does
alone was never asked until it was asked directly
(Sec.~\ref{sec:also-preservation}).

\subsection{Single-Term Ablations}
\label{sec:also-single}

Once the corpus was widened and the trust region removed, each candidate term
could be added individually to a common baseline.
Table~\ref{tab:ablation} reports arms trained from released M7 under the
optimizer of Sec.~\ref{sec:crossres}, all declaring the same horizon and all
terminated at the same milestone, differing only in the named flag.

\begin{table}[t]
\caption{Single-term ablations on the wide corpus with no trust region, best
frozen milestone, macro-scroll Dice. The seed replicate differs from the
baseline only in its random seed and bounds what any single arm can claim.}
\label{tab:ablation}
\centering
\small
\begin{tabular}{lr}
\toprule
Arm & Macro Dice \\
\midrule
Occupancy only (CE + Dice) & $\ablBaseline$ \\
Occupancy only, seed replicate & $\ablSeedReplicate$ \\
\midrule
$+$ separation shell $0.5$ & $\ablSeparation$ \\
$+$ M7 function anchor $0.5$ & $\mathbf{\ablKL}$ \\
\bottomrule
\end{tabular}
\end{table}

The anchor is the one term that clears its own noise floor, and it is the one
term in Eq.~\ref{eq:objective}.
The separation arm lands below the baseline.
The two occupancy-only arms differ from each other by more than most effects we
were trying to detect, which is why every later comparison is expressed against
a measured seed band rather than against zero.

\subsection{Separation Shell}
\label{sec:also-separation}

The shell placed background cross-entropy on a two-voxel dilation of the
teacher-supported positives, intersected with confidently empty teacher space.
The intent was a local fence: attack unsupported girth and short cross-wrap
welds while leaving the occupancy term free to grow faint material elsewhere.

Measured as a dial, it does what it claims.
On a separate release-ladder sweep, raising its weight from $2$ to $3$ cut
reference-relative bridge pixels by roughly half.
It also cost about sixteen points of reference skeleton recall, against a
tolerance of one, and lowering it toward zero inverted both effects.
What the sweep does not show is any movement of the underlying trade.
Compared at matched coverage rather than at a fixed threshold, a model trained
with the shell sits on the same bridge-versus-coverage curve as one trained
without it; the shell selects a point on that curve, and so does the
segmentation threshold, which is free.
We retired it because a training-time term that duplicates a deployment-time
dial should be the dial.

\subsection{Medial Crest Recall}
\label{sec:also-crest}

The crest term supervised sheet \emph{existence} separately from amount of
material.
The fine teacher's medial surface was extracted with the center--radius view of
LSMAT~\cite{rebain2019lsmat}, projected to the coarse grid with a max operator
so that a thin center survives downsampling, and rewarded through a per-sample
recall term evaluated only on crest voxels.
Slice-wise radii were deliberate: pairing 2D centers with a 3D distance
transform would substitute spheres for disks and collapse the medial surface
toward a curve.

Its premise was the occupancy ceiling of Sec.~\ref{sec:crossres}.
That premise is corpus-dependent, and on the wide corpus the diagnostic that
was supposed to trigger the term does not fire: the fraction of
continuation-band validation voxels already above the lowest gate-passing
threshold is comfortably above the level at which a ceiling would be indicated.
The term was accordingly never triggered for a wide-corpus arm of its own.
Where it was measured, on the narrower release ladder, raising its weight
reduced bridge burden by roughly a fifth but cost five points of skeleton
recall against the same one-point tolerance, and the rendered panels showed it
refilling broad mass rather than extending thin lines.
Global recall pressure turned out to be the wrong shape for the problem: it
buys length by buying width.

A focused variant survives as untested rather than refuted.
Instead of rewarding all crest voxels, it applies a one-sided hinge to the
weakest fifth of them, so gradient vanishes once a selected voxel clears a floor.
It had the best burden-versus-recall behaviour of the medial family in its one
trial, and it cannot be evaluated on the present corpus at all: the term
requires a medial sidecar on every training row, and the human-labelled rows of
the wide corpus can never carry one.
Testing it means building a crest-complete corpus first, which changes the
corpus and therefore needs its own control arm.

\subsection{M7 Positive Preservation}
\label{sec:also-preservation}

Preservation applied a one-sided foreground cross-entropy to voxels where the
frozen M7 predicted foreground inside a two-voxel corridor around teacher
positives.
Its stated job was to stop de-blobbing from deleting M7-supported sheet, and it
had the best surviving motivation of any retired term: our students beat M7 on
overlap and on merge behaviour but lose to it on topology on the one
human-labelled holdout, which is the deficit preservation is shaped to fix.

It was also, of the seven, the only term never measured in isolation at scale,
so we measured it.
Three arms were preregistered against the shipped recipe, each declaring the
full horizon and terminating at a common milestone: the original hard form at
its shipped weight; a one-sided soft-floor form that resists only where the
student has fallen below M7; and, conditional on a rim-thickening trigger, the
hard form paired with a reduced separation shell.
Kill criteria, a seed-derived tie margin, and the thickening trigger were fixed
before any arm ran.

All three were rejected, and the failure mode was the one the preregistration
named.
Pooled paired deltas against the shipped recipe over five matched milestones
were $\presHard$ for the hard form, $\presSoft$ for the soft floor, and
$\presSep$ for the shell-paired arm, every interval excluding zero.
Each arm failed the blind six-cube morphology gate at two consecutive
milestones on interior fraction and maximum thickness.
Foreground ran between $1.3$ and $2.3$ times the reference where the shipped
recipe runs at $0.71$, and reference-relative bridge pixels ran an order of
magnitude above it.
The frontier veto could not even be evaluated: at the coverage these models
reach, they lie beyond the end of the reference curve, which is itself the
finding.
None of the arms was inert---the term contributed between 13\% and 22\% of the
committed training loss---so these are measured negatives rather than absences.

The shell-paired arm is the informative one.
It halves the damage on every axis and still fails every gate.
Read together, the three arms say that preservation's cost is paid in girth and
that the separation shell of the original stack existed to pay it.
Reconstructing that pair means reintroducing a term we retired in
Sec.~\ref{sec:also-separation} for duplicating a free deployment-time dial.

\subsection{Dynamic Connectivity}
\label{sec:also-connectivity}

The connectivity term asked, existentially, for some viable route through an
audited teacher-medial corridor joining two or more disconnected M7 components.
Its precursor prescribed one minimum-off-axis route and raised the weakest
tenth of its logits; that version tied the incumbent at the smallest weight and
degraded ordinary validation monotonically at larger ones, which is what
motivated replacing a prescribed path with a bottleneck on whichever path the
student currently supports.
The mechanism is sound and we still regard it as the most interesting of the
retired terms.

It is nonetheless inapplicable to the corpus that produced our result, for a
structural reason rather than a measured one.
An event requires a full-volume M7 prediction to define the component-contact
pins, and that prediction exists for exactly one of the training scrolls.
The frozen atlas contains 149 events, all from that scroll; on the wide corpus
roughly three percent of rows carry any event at all, so the term is close to
inert by construction and its weight cannot be interpreted.
Producing the missing full-volume predictions is days of inference, and until
that exists the honest status is untested at breadth, not refuted.

\begin{figure*}[t]
  \centering
  \figureasset{figure_crossres_supervision}{0.96\linewidth}{2.0in}
  \caption{Supervision fields from the retired seven-term configuration, shown
  on actual training data at dynamic event 138. Projected soft occupancy is
  retained by the present method; the sparse medial crest, the radius-two
  separation shell, the teacher-medial corridor, the component-contact pins and
  the free anchors supplied to the max--min connectivity recurrence all belong
  to terms this appendix retires.}
  \Description{Training slice with projected teacher occupancy, medial crest,
  de-blob shell, M7 fragments, and dynamic-connectivity metadata.}
  \label{fig:supervision}
\end{figure*}

\subsection{Terms Proposed and Rejected Later}
\label{sec:also-later}

Two further candidates were measured after the reduction and are recorded for
completeness.
A hardened Dice term, taking the thresholded label the benchmark actually
scores in addition to the soft target, lost to the shipped recipe at both
matched milestones with intervals excluding zero, and carries a documented
confound: human rows are already binary, so their Dice weight doubles.
Dropping the human-labelled rows entirely lost at the earlier milestone and won
narrowly at the later one, inside the seed band, and was not promoted.

\subsection{The Retired Configuration, for Reference}
\label{sec:also-v31}

Table~\ref{tab:duration} preserves the completed duration ladder of the
seven-term configuration on its single-scroll corpus, as originally reported.
Its best calibrated macro-scroll Dice was $\bestMacro$ at \bestSamplesText\
samples, a gain of roughly $0.02$ over released M7.
The present method gains $\fzGain$ on the frozen benchmark.
The entire seven-term stack was worth about a sixth of what corpus breadth and
a completed schedule were worth, which is the ratio we would most like a reader
to take away.

\begin{table}[t]
\caption{Historical duration ladder of the retired seven-term configuration.
``Cal.'' is calibrated macro-scroll Dice; ``$t=.40$'' uses a fixed threshold;
``gain'' is calibrated macro gain over the M7 initial model. All calibrated
thresholds equal the $0.25$ search boundary, so the optimum is censored.}
\label{tab:duration}
\centering
\small
\begin{tabular}{rrrrr}
\toprule
Samples & Cal. & $t=.40$ & Gain & Min. gain \\
\midrule
\durationRows
\bottomrule
\end{tabular}
\end{table}
